\chapter{Tổng kết}
\label{ch:tongket}

\section{Kết quả đạt được}

Tính đến bản cập nhật này, luận văn đã: (1)~phát biểu lại bài toán theo mô hình
triển khai thực tế (local/on-device, granularity theo mục đích, $(\alpha,\delta)$-hữu
dụng, bảo vệ endpoint) với bảng yêu cầu ánh xạ vào harm thực tế; (2)~chỉ ra và
chứng minh hình thức các lỗi phá vỡ đảm bảo Geo-I trong pipeline Thực tập 2;
(3)~đề xuất REM --- cơ chế lũy thừa trên \emph{toàn bộ} tập đỉnh lưới đường cố
định với chứng minh $\eps$-Geo-I Euclid cho ideal kernel (bản Euclid của
GEM~\cite{takagi2019geo}), đầu ra nằm trên đường theo cấu trúc; (4)~đề xuất T-REM
--- trọng số khả đạt chỉ phụ thuộc thông tin công khai, không đổi đảm bảo
per-release, \emph{giảm} tín hiệu speed-implausibility (không claim đóng hẳn
velocity attack); (5)~đề xuất SM-REM (đảm bảo cho vị trí tĩnh lặp lại,
Định lý~\ref{thm:smrem}) và \textbf{PR-SM-REM} (reuse riêng tư, chặn per-release
$(\eps_{\text{test}}{+}\eps_{\text{release}})$-Geo-I hữu hạn,
Định lý~\ref{thm:prsmrem}) --- nhắm kịch bản averaging mà bản PTPPM
stateless~\cite{min2025road} bỏ ngỏ; (6)~bộ khung đánh giá theo chuẩn văn liệu
(Bayes/HMM proxy dạng-REM, averaging đa-home 40$\times$8 với bootstrap CI,
quality loss, k-NN POI tổng hợp, tính thực tế) trên dữ liệu GPS thật; (7)~hệ
thống mô phỏng ba góc nhìn gắn với kịch bản đời thực. Mọi đảm bảo là cho ideal
real-arithmetic kernel (Nhận xét~\ref{rmk:finiteprec}).

\section{Hạn chế hiện tại}

\textbf{Hiện thực hữu-hạn-độ-chính-xác không phải pure Geo-I.} Sampler Gumbel-max
và phép thử Laplace dùng RNG 53-bit nên có zero-support phụ thuộc input
(Nhận xét~\ref{rmk:finiteprec}); đảm bảo chỉ đúng cho ideal kernel. Exact discrete
sampler = hướng phát triển.

\textbf{Attacker là proxy, chưa optimal.} Bayes/HMM/averaging dùng likelihood
dạng-REM, chính xác cho REM nhưng là proxy cho T-REM/SM-REM/PR-SM-REM; GPS jitter
chưa vào likelihood. Một attacker sequential đúng của PR mạnh hơn proxy (chẩn đoán
độc lập của verifier). Do đó \emph{không} dùng chênh lệch proxy để claim cơ chế
này vượt cơ chế kia; cần attacker mechanism-aware + jitter-aware đầy đủ.

\textbf{Đảm bảo trajectory chưa đầy đủ.} PR-SM-REM mới có chặn per-release; một
scalar $\eps_w$ w-event cần budget manager theo cửa sổ (chưa cài). Chống-averaging
của PR-SM-REM là \emph{hữu-hạn-horizon}: cả nhánh reuse lẫn resample đều ngẫu nhiên
mỗi bước (tại $d{=}\theta$ xác suất resample $=1/2$), nên không tuyệt đối và khác
tính chất static-same-cell của SM-REM. SM-REM chỉ đảm bảo vị trí tĩnh lặp lại (một
ô, cache khởi tạo rỗng); pattern revisit rò rỉ.

\textbf{Phạm vi \& provenance.} Artifact là reference mechanism + simulator Flask
phía server, chưa phải hiện thực on-device (không có latency đo cách ly);
cache/định danh lifecycle (R8--R12) ngoài phạm vi. Đã bổ sung: build script tất
định + kiểm hash graph fail-closed, seed dẫn xuất từ khóa ngữ nghĩa (permutation
test), provenance \texttt{msc-experiment-v1} (source commit + dirty flag, record
IDs, raw rows), và dependency lock dev (\texttt{requirements-dev.txt}). Còn lại:
moving benchmark một-seed (exploratory, chưa multi-seed CI), averaging dùng
stay-point (không phải nhãn cư trú), clean-clone vẫn cần raw data ngoài Git. Chưa
chạy RAoPT; chưa có comparator thực thi cho GEM (metric đường đi ngắn nhất)/PTPPM/
Eclipse (mới ở mức prose). Chi tiết đối chiếu ở \texttt{docs/reviews/}.

\section{Hướng phát triển}

\begin{enumerate}
  \item Chạy tấn công RAoPT (mã nguồn công khai) đối kháng cả bốn cơ chế ---
  chuẩn đánh giá tối thiểu theo SoK PoPETs 2024~\cite{buchholz2024sok};
  \item Hiện thực GEM với khoảng cách đường đi ngắn nhất và so sánh REM--GEM
  theo giao thức của~\cite{takagi2019geo};
  \item Phiên bản $\delta$-location set trên lưới đường: tập ứng viên là các
  đỉnh khả đạt từ vùng đã công bố, cho đảm bảo kiểu PIM~\cite{xiao2015protecting}
  dưới kẻ tấn công Markov;
  \item Mở rộng thực nghiệm: Porto Taxi (thành phố thứ hai), 100+ quỹ đạo,
  nhiều seed với khoảng tin cậy, đường cong Pareto riêng tư--tiện ích dày hơn;
  \item Metric co giãn theo ngữ nghĩa trên đồ thị đường (``elastic GG-I'') ---
  khoảng trống chưa có công bố chiếm giữ, kết hợp
  \cite{chatzikokolakis2015elastic} với \cite{takagi2019geo}.
\end{enumerate}
